\documentclass{article}

% custom margins
\usepackage[a4paper,margin=1.5in]{geometry}
\renewcommand{\baselinestretch}{1.2}
%\AddToHook{cmd/section/before}{\clearpage}

% emma's long list of custom macros and universally used packages
\include{../macros-and-packages.tex}
\newcommand{\seperator}{\vspace{-8pt}\hline\vspace{8pt}}

% colored box behind proofs
\tcolorboxenvironment{proof}{
	colback=white,
	boxrule=0pt,
	frame hidden,
	borderline west={1pt}{0pt}{black},
	before skip=0.75cm,
	after skip=0.75cm,
	sharp corners,
	breakable,
	enhanced,
}

\renewcommand*\contentsname{Inhalt}
\renewcommand*\proofname{Lösung}

\newcommand{\ex}[1]{\tit{Übung} #1.}
\pagestyle{fancy} %allows headers

\lhead{Emma Bach}
\rhead{\today}
\title{\bfseries Aufgabensammlung Algebra und Zahlentheorie}
\date{Wintersemester 25/26}
\author{Emma Bach}
\begin{document}
	\maketitle
	\section*{\centering - Gruppen -}
		\ex{0.2} Man gebe ein unverkürzbares Erzeugendensystem der Gruppe $\bZ^2$ an, welches aus drei Elementen besteht.
		\newpar
		\ex{1.2} Wie viele Untergruppen hat die additive Gruppe eines dreidimensionalen Vektorraums über dem Körper mit drei Elementen?
		\begin{proof}
			\begin{enumerate}
				\item Das neutrale Element erzeugt die triviale Untergruppe.
				\item Jedes der 26 anderen Elemente erzeugt eine Untergruppe, welche zu $\bZ/3\bZ$ isomorph ist. Dabei erzeugen $v$ und $-v$ die selbe Untergruppe, es gibt also $13$ "Klassen" von linear unabhängigen Vektoren.
				\item Jedes Paar von linear unabhängigen Vektoren erzeugt ein Untergruppe, welche zu $\bZ/3\bZ \times \bZ/3\bZ$ isomorph ist. Da ein solches Paar dadurch definiert werden kann, welche linear unabhängige Klasse nicht enthalten ist, gibt es hiervon ebenfalls $13$.
				\item Zuletzt ist $\bZ/3\bZ \times \bZ/3\bZ \times \bZ/3\bZ$ selbst eine Untergruppe.
			\end{enumerate}
			Insgesamt gibt es also $28$ Untergruppen.
		\end{proof}
		\newpar
		\ex{1.2.B} Wie viele Untergruppen hat allgemein $(\bZ/q\bZ)^n$?
		\newpar
		\ex{1.3} Man bestimme alle Untergruppen von $\bZ/20\bZ$.
		\begin{proof}
			Die Untergruppen sind $d\bZ/20\bZ \cong (\bZ/\frac{20}{d}\bZ)$, für jeden Teiler $d$ von $20$.
		\end{proof}
		\newpar
		\ex{1.4} Sind $H, K \subset G$ zwei Untergruppen einer endlichen Gruppe $G$ mit $\abs{H}$ teilerfremd zu $\abs{K}$ und $\abs{H} \cdot \abs{K} = \abs{G}$, so induziert die Verknüpfung eine Bijektion $H \times K \to G$.
		\begin{proof}
			Betrachte die Multiplikationsabbildung $m : H \times K \to G$, $(h,k) \mapsto hk$. Wir wollen zeigen, dass diese injektiv ist, dann folgt aus $\abs{H \times K} = \abs{H} \cdot \abs{K} = \abs{G}$ gemäß Schubfachprinzip auch Bijektivität.
			\newpar
			Sei also $h_1k_1 = h_2k_2$. Durch Multiplikation mit $h_2^{-1}$ von Links und mit $k_1^{-1}$ von Rechts folgt
			\begin{align*}
				h_2^{-1}h_1 
				&= h_2^{-1} h_1 k_1 k_1^{-1}\\
				&= h_2^{-1} h_2 k_2 k_1^{-1}\\
				&= k_2k_1^{-1},
			\end{align*}
			also insbesondere
			\begin{align*}
				h_2^{-1}h_1 = k_2k_1^{-1} \in H \cap K
			\end{align*}
			$H \cap K$ ist sowohl von $H$ als auch von $K$ eine Untergruppe. Da $\abs{H}$ und $\abs{K}$ teilerfremd sind, muss $H \cap K$ also die triviale Gruppe sein, also gilt insbesondere $h_2^{-1}h_1 = k_2k_1^{-1} = 1$, woraus $h_1 = h_2$ und $k_1 = k_2$ folgt. Somit ist die Abbildung injektiv, also bijektiv.
		\end{proof}
		\newpar
		\ex{2.1.a)} Gegeben zwei Elemente $x,y$ endlicher Ordnung einer abelschen Gruppe $G$ teilt die Ordung von $xy$ das kleinste gemeinsame Vielfache der Ordnungen von $x,y$.
		\begin{proof}
			Es gilt $g^k = e$ gdw. $\ord(g) \mid k$. Sei also $n = \kgV(\ord x, \ord y)$. Dann folgt direkt $x^n = e = y^n$. Da $G$ abelsch ist, folgt
			\begin{align*}
				(xy)^n = x^ny^n = e,
			\end{align*}
			also $\ord(xy) \mid \kgV(\ord x, \ord y)$.
		\end{proof}
		\newpar
		\ex{2.1.b)} Sind die Ordnungen von $x$ und $y$ Teilerfremd, so gilt sogar
		\begin{align*}
			\ord(xy) = (\ord x) \cdot (\ord y)
		\end{align*}
		\begin{proof}
			Da $\ggT(\ord x, \ord y) = 1$ folgt 
			\begin{align*}
				\ord x \cdot \ord y = \kgV(\ord x, \ord y).
			\end{align*} 
			Wir wissen bereits $\ord(xy) \mid \kgV(\ord x, \ord y)$, also folgt $\ord(xy) = \kgV(\ord x, \ord y)$ bereits, falls wir 
			\begin{align*}
				\kgV(\ord x, \ord y) = \ord x \cdot \ord y \mid \ord(xy)
			\end{align*} zeigen können.
			\newpar
			Es gilt per per Definition 
			\begin{align*}
				(xy)^{\ord xy} = e,
			\end{align*}
			also folgt aus der Kommutativität der Gruppenoperation wieder
			\begin{align*}
				x^{\ord xy} y^{\ord xy} = e,
			\end{align*}
			also insbesondere
			\begin{align*}
				x^{\ord xy} = y^{-{\ord xy}}
			\end{align*}
			Es gilt also $x^{\ord xy} \in \angles{y}$ und $y^{-{\ord xy}} \in \angles{x}$, also wie in 1.4 insbesondere 
			\begin{align*}
				x^{\ord xy} \in \angles{y} \cap \angles{x}.
			\end{align*}
			Da die Ordnungen von $x$ und $y$ teilerfremd sind folgt wieder durch Lagrange $\angles{y} \cap \angles{x} = \angles{e}$. Also gilt $x^{\ord xy} = e$ und $y^{\ord xy} = (y^{-{\ord xy}})^{-1} = e^{-1} = e$, also sind $\ord x$ und $\ord y$ beides Teiler von $\ord xy$, also folgt aus der Teilerfremdheit der beiden Ordnungen
			\begin{align*}
				\kgV(\ord x, \ord y) = \ord x \cdot \ord y \mid \ord(xy),
			\end{align*}
			was zu zeigen war.
		\end{proof}
		\newpar
		\ex{2.3} Wie viele Gruppenhomomorphismen gibt es von $\bZ/8\bZ$ zur symmetrischen Gruppe $\cS_4$?
		\begin{proof}
			Gemäß universeller Eigenschaft des Quotienten gilt
			\begin{align*}
				\Hom(\bZ/8\bZ, \cS_4) = \lr{f \in \Hom(\bZ, \cS_4) \mid f(8\bZ) = \id}
			\end{align*}
			Gemäß Vorlesung gilt $\Grp(\bZ, G) \cong G$, also
			\begin{align*}
				\lr{f \in \Hom(\bZ, \cS_4) \mid f(8\bZ) = \id}
				= \lr{\sigma \in \cS_4 \mid \sigma^8 = \id}
			\end{align*}
			Wir suchen also alle Elemente von $\cS_4$, deren Ordnung $8$ teilt. Die Elemente von $\cS_4$ sind:
			\begin{enumerate}
				\item Die Identität (Ordnung $1$),
				\item Sechs Transpositionen (Ordnung $2$),
				\item Drei Doppeltransmutationen (Ordnung $2$),
				\item Acht $3$-Zyklen (Ordnung $3$),
				\item Sechs $4$-Zyklen (Ordnung $4$).
			\end{enumerate}
			Es kommen also alle Elemente bis auf die $3$-Zyklen in Frage, insgesamt gibt es also $24 - 8 = 16$ Gruppenhomomorphismen.
		\end{proof}
		\newpar
		\ex{2.4} Gegeben einen surjektiven Gruppenhomomorphismus $\phi : G \surj \ol G$ und einen Normalteiler $\ol N \subset \ol G$ mit Urbild $\phi^{-1}(\ol N) := N \subset G$ induziert $\phi$ einen Gruppenisomorphismus
		\begin{align*}
			\phi_q : G/N \iso \ol G / \ol N
		\end{align*}
		\begin{proof}
			Sei $q : \ol G \to \ol G / \ol N$ die kanonische Quotientenabbildung. Wir betrachten die Verktettung $\phi_q := q \circ \phi$. Da $q$ und $\phi$ beide surjektiv sind, ist $\phi_q$ ebenfalls surjektiv. Der Kern-Bild-Isomorphiesatz gibt uns nun
			\begin{align*}
				\phi_q : G/\ker \phi_q \iso \im~ \phi_q = \ol G / \ol N
			\end{align*}
			Und für den Kern gilt:
			\begin{align*}
				\ker(\phi_q) 
				&= (q \cdot \phi)^{-1}([e_{\ol G}])\\
				&= \phi^{-1} q^{-1} ([e_{\ol G}])\\
				&= \phi^{-1}(\ol N)\\
				&= N.
			\end{align*}
			Also haben wir wie gewünscht
			\begin{align*}
				\phi_q : G/N \im~ \phi_q = \ol G / \ol N.
			\end{align*}
		\end{proof}
		\newpar
		\ex{3.2.a} Jede Untergruppe einer zyklischen Gruppe ist zyklisch.
		\begin{proof}
			Sei $(G,+)$ eine zyklische Gruppe mit Erzeuger $a \in G$. Sei $H$ eine Untergruppe von $G$. Wäre $H$ nicht zyklisch, so gäbe es Elemente $h, h' \in H$, sodass kein Element von $H$ sie beide erzeugt.
			\newpar
			Da $G$ selbst zyklisch ist, existieren $n, n' \in \bZ$, sodass $h = n \cdot a$ und $h' = n' = a$. Durch den Euklidischen Algorithmus finden wir $r,s \in \bZ$, sodass
			\begin{align*}
				\kgV(n, n') = rn + sn'
			\end{align*}
			Nun gilt allerdings:
			\begin{align*}
				\kgV(n, n') \cdot a
				&= (rn + sn') \cdot a\\
				&= rna + sn'a\\
				&= rh + sh' \in H,
			\end{align*}
			und da $\kgV(n, n')$ sowohl $n$ als auch $n'$ teilt, erzeugt $\kgV(n, n') \cdot a \in H$ sowohl $h$ als auch $h'$.
		\end{proof}
		\newpar
		\ex{3.2.b} Für beliebiges $m \in \bN$ hat jede Untergruppe von $\bZ/m\bZ$ die Form $d\bZ/m\bZ$ für einen Teiler $d$ von $m$, und jeder Teiler liefert eine solche Untergruppe.
		\begin{proof}
			Wir haben soeben gezeigt, dass jede Untergruppe von $\bZ/m\bZ$ zyklisch ist. Da nach Lagrange die Ordnung jeder Untergruppe die Ordnung der Gruppe teilt, muss sie also auch von einem Element der entsprechenden Ordnung erzeugt werden.
			\newpar
			Jedes Element $k \in \bZ/m\bZ$ hat Ordnung
			\begin{align*}
				\ord k = \frac{m}{\ggT(k,m)}
			\end{align*} 
			Sei also $H$ eine Untergruppe von $\bZ/m\bZ$ der Ordnung $\frac{m}{d}$. Gemäß der eben genannten Gleichung muss jeder Erzeuger $k$ von $H$ $d$ als Teiler haben. Somit muss jedes solches $k$ in der von $d$ erzeugten Untergruppe $d\bZ/m\bZ$ liegen, welche ebenfalls Ordnung $\frac{m}{d}$ hat. Somit erzeugt jedes solches $k$ die selbe Untergruppe. Gleichzeitig gibt es für jedes $d$ nur eine solche Untergruppe, also stehen die gefragten Mengen in Bijektion.
		\end{proof}
		\newpar
		\ex{3.3} 
		\begin{enumerate}
			\item Wie viele Erzeuger hat die multiplikative Gruppe eines Körpers mit elf Elementen?
			\item Man gebe einen solchen Erzeuger an.
			\item Wie viele Elemente des Körpers mit elf Elementen sind Quadrate?
			\item Wie viele Elemente des Körpers mit elf Elementen sind dritte Potenzen?
		\end{enumerate}
		\begin{proof}
			Durch Ausprobieren trivial :)
		\end{proof}
		\newpar
		\ex{4.1} Seien $H \supset N$ eine Gruppe mit einem Normalteiler und $X$ eine Menge mit $H$-Operation. Man zeige, dass es auf dem Bahnenraum $X/N$ genau eine Operation der Quotientengruppe $H/N$ gibt, sodass
		\begin{align*}
			(hN)(Nx) = Nhx
		\end{align*}
		\begin{proof}
			\begin{enumerate}
				\item 
					Wir wollen zuerst zeigen, dass aus $Nx = Ny$ für jedes $h \in H$ auch $Nhx = Nhy$ folgt.
					\newpar
					Da $Nx = Ny$ existiert ein $n \in N$, sodass $x = ny$. Es folgt $hx = hny$. Da $N$ ein Normalteiler ist, gilt desweiteren  
					\begin{align*}
						n' := hnh^{-1} \in N.
					\end{align*}
					Für dieses $n'$ folgt $hn = n'h$, also $hx = n'hy$. Da $n' \in N$ folgt also
					\begin{align*}
						Nhx = Nhy
					\end{align*}
				\item Als nächstes wollen wir aus dieser Vorüberlegung die Wohldefiniertheit der definierten Operation folgern, also zeigen, dass zusätzlich für $g,h \in H$ aus $gN = hN$ auch $(gN)(Nx) = (hN)(Nx)$ folgt.
				\newpar
				Da $N$ ein Normalteiler ist, gilt auch $gN = Nh$, also gibt es ein $n \in N$, sodass $g = nh$, also $gx = nhx$, also $Ngx = Nhx$. Somit gilt
				\begin{align*}
					(gN)(Nx) = Ngx = Nhx = (hN)(Nx).
				\end{align*}
				\item Zuletzt wollen wir zeigen, dass diese Verknüpfung tatsächlich eine Gruppenoperation definiert. Zu zeigen ist also:
				\begin{enumerate}
					\item $(1N)(Nx) = Nx$. Dies gilt trivial, da $(1N)(Nx) = N1x = Nx$.
					\item Sei $g,h \in G$. Dann wollen wir 
					\begin{align*}
						((gN)(hN))(Nx) = (gN)((hN)(Nx))
					\end{align*}
					Dies folgt ebenfalls recht schnell:
					\begin{align*}
						((gN)(hN))(Nx) 
						&= (ghN)(Nx)\\
						&= N(gh)x\\
						&= N(g)(hx)\\
						&= (gN)(Nhx)\\
						&= (gN)((hN)(Nx)
					\end{align*}
				\end{enumerate}
			\end{enumerate}
		\end{proof}
		\newpar
		\ex{4.3} Man gebe ein Repräsentantensystem an für die Konjugationsklassen der Gruppe aller Isometrien der euklidischen Ebene.
		\newpar
		\ex{4.4} Man bestimme alle Kompositionsreihen der Gruppe der Drehsymmetrien eines Würfels.
		\newpar
		\ex{5.1} Sei $G$ eine endliche Gruppe mit einem Normalteiler $N$ und $p$ prim. Dann ist eine Untergruppe $P \subset G$ genau dann eine $p$-Sylow von $G$, wenn $P \cap N$ eine $p$-Sylow von $N$ ist und das Bild von $P$ in $G/N$ eine $p$-Sylow von $G/N$.
		\begin{proof}\phantom{}
			\begin{enumerate}
				\item Sei $P$ eine $p$-Sylow von $G$.
				\begin{enumerate}
					\item Da $P$ eine $p$-Sylow ist, gilt $p \nmid \frac{\abs{G}}{\abs{P}}$. Da $NP$ eine Untergruppe von $G$ ist, gilt $\frac{\abs{NP}}{\abs{P}} \mid \frac{\abs{G}}{\abs{P}}$, also insbesondere $p \nmid \frac{\abs{NP}}{\abs{P}}$.
					\newpar
					Wenden wir den Isomorphiesatz auf $f : P \to NP/N$, $p \mapsto Np$ an, erhalten wir
					\begin{align*}
						P/(P \cap N) = P/(\ker f) \cong \im f = NP/N
					\end{align*}
					somit gilt
					\begin{align*}
						p \nmid \frac{\abs{NP}}{\abs{P}} = \frac{\abs{N}}{\abs{N \cap P}}
					\end{align*}
					Da $N \cap P$ eine $p$-Gruppe in $N$ ist, folgt daraus auch, dass $N \cap P$ eine $p$-Sylow von $N$ ist.
					\item (...)
				\end{enumerate}
				\item Sei $N \cap P$ eine $p$-Sylow von $N$ und $q(P) = P/N$ eine $p$-Sylow von $G/N$. Dann ist $P$ "offensichtlich" eine $p$-Gruppe in $G$. Es bleibt noch zu zeigen, dass $p \nmid \frac{\abs{G}}{\abs{P}} = [G : P]$. Durch den Isomorphiesatz folgt weiterhin
				\begin{align*}
					\frac{\abs{NP}}{\abs{P}} = \frac{\abs{N}}{\abs{N \cap P}}
				\end{align*}
				Desweiteren gilt
				\begin{align*}
					[G : NP] = [G/N : P/N]
				\end{align*}
				Es gilt nun
				\begin{align*}
					[G : P] &= [G : NP][NP : P]\\
					&= [G/N : P/N][N : N \cap P]
				\end{align*}
				Da $P/N$ eine $p$-Sylow von $G/N$ und $N \cap P$ eine $p$-Sylow von $N$ ist, teilt $p$ keinen dieser beiden Terme, also auch nicht $[G:P]$, was zu zeigen war.
			\end{enumerate}
		\end{proof}
		\newpar
		\ex{5.2} Für jede Primzahl $p$ sind die einzigen Gruppen der Ordnung $2p$ bis auf Isomorphismus die zyklische Gruppe $\bZ/2p\bZ$ und die Diedergruppe $D_p$.
		\begin{proof}
			Sei $G$ eine Gruppe der Ordnung $2p$. Gemäß viertem Sylowsatz ist die Zahl $n_2$ der $2$-Sylows von $G$ ein Teiler von $\frac{\abs{\bZ/2p\bZ}}{2} = p$, also entweder $1$ oder $p$. Das zweite Kriterium aus dem vierten Sylowsatz besagt $n_2 \cong 1 \mod 2$, was keinen der beiden Fälle ausschließt.
			\newpar
			Desweiteren ist die Anzahl $n_p$ der $p$-Sylows von $G$ ein Teiler von $\frac{\abs{\bZ/2p\bZ}}{2} = 2$, also entweder $1$ oder $2$. In diesem Fall besagt das zweite Kriterium $n_p \cong 1 \mod p$, was den Fall $n_p = 2$ ausschließt.
			\newpar
			Der Fall $p = 2$ ist bereits bekannt, sei also $p > 2$ angenommen.
			\begin{enumerate}
				\item Angenommen, $n_2 = p$. Jede dieser Gruppen hat per Definition $2$ Elemente, ist also isomorph zur zyklischen Gruppe $\bZ/2\bZ$. Somit hat $G$ $p$ Elemente der Ordnung $2$. Sei $a$ ein solches Element. Gemäß Satz von Cauchy existiert desweiteren ein Element $b$ der Ordnung $p$.
				\newpar
				Gemäß Satz von Lagrange gilt $\ord ab \in \lr{1,2,p,2p}$.
				\begin{enumerate}
					\item Im Fall $(\ord ab) = 2p$ wäre $G$ zyklisch, was der Annahme wiedersprechen würde, dass $G$ $p$ $2$-Sylows hat.
					\item Im Fall $(\ord ab) = 1$ hätten wir $p = 2$, (?) was wir ausgeschlossen hatten. 
					\item Im Fall $(\ord ab) = p$ wäre $\angles{ab}$ eine zweite $p$-Sylow, was wir ausgeschlossen hatten.
				\end{enumerate} 
				Somit gilt $(\ord ab) = 2$. Wir erhalten die Präsentation
				\begin{align*}
					G = \angles{a,b \mid \ord b = p, \ord a = \ord ab = 2},
				\end{align*}
				welche eindeutig die Diedergruppe $D_p$ definiert.
				\item Angenommen, $n_2 = 1$. Da $G$ auch nur eine $p$-Sylow hat, sind beide Sylowuntergruppen normal. Es folgt, dass der Schnitt der beiden Gruppen die triviale Gruppe ist, also ist $G$ das direkte Produkt der beiden, also $G \cong \bZ/2\bZ \times \bZ/p\bZ$. Gemäß Chinesischem Restsatz ist nun $\bZ/2\bZ \times \bZ/p\bZ$ isomorph zu $\bZ/2p\bZ$, also $G$ zyklisch.
			\end{enumerate}
		\end{proof}
		\newpar
		\ex{5.3} Sei $p > q$ und $q$ kein Teiler von $p-1$. So ist jede Gruppe der Ordnung $pq$ zyklisch.
		\begin{proof}
			Sei $G$ eine Gruppe der Ordnung $pq$. Gemäß Lagrange hat jedes Element von $G$ die Ordnung $1,p,q,$ oder $pq$. Insbesondere erzeugen die Elemente der Ordnung $p$ und $q$ die jeweiligen Sylowgruppen.
			\newpar
			Gemäß der Sylowsätze ist die Anzahl $n_p$ der $p$-Sylows von $G$ ein Teiler von $q$, also $n_p = 1$ oder $n_p = q$. Gleichzeitig muss $n_p \cong 1 \mod \frac{\abs{G}}{q} = p$, also $p \mid n_p - 1$. Da $p > q$ und beide Zahlen prim sind, ist dies im Fall $n_p > 1$ unmöglich, also folgt $n_p = 1$. Diese $p$-Sylow enthält $p$ Elemente, von denen $p-1$ Erzeuger sind. Somit gibt es $p-1$ Elemente der Ordnung $p$.
			\newpar
			Analog ist die Anzahl $n_q$ der $q$-Sylows von $G$ ein Teiler von $p$, der Kongruent zu $1$ modulo $q$ ist. Wäre $n_q = p$, hätten wir $p \cong 1 \mod q$, also $q \mid p-1$, was wir ausgeschlossen hatten. Also existiert auch genau eine $q$-Sylow, also existieren $q-1$ Elemente der Ordnung $q$.
			\newpar
			Zusammen mit dem Neutralen Element erhalten wir $p + q - 1$ Elemente der Ordnung ungleich $pq$. Für Primzahlen $p > q$ gilt $p + q - 1 < pq$, also ist mindestens ein Element übrig, welches demenstprechend Ordnung $pq$ hat und somit die Gruppe $G$ erzeugt. Somit ist $G$ zyklisch.
		\end{proof}
	\clearpage
	\section*{\centering - Ringe -}
		\ex{5.4} In jedem Ring $R$ gilt für alle $a,b \in R$:
		\begin{enumerate}
			\item $0 \cdot a = 0$
			\item $-a = (-1) \cdot a$
			\item $(-1) \cdot (-1) = 1$
			\item $(-a) \cdot (-b) = ab$
		\end{enumerate}
		\begin{proof}
			\begin{enumerate}
				\item 
				\begin{align*}
					0 
					&= 0a + (-0a)\\ 
					&= (0 + 0)a + (-0a)\\
					&= 0a + 0a + (-0a)\\
					&= 0a
				\end{align*}
				\item
				\begin{align*}
					a + (-1) \cdot a 
					&= 1 \cdot a + (-1) \cdot a\\
					&= (1 + (-1)) \cdot a\\
					&= 0a\\
					&= 0
				\end{align*}
				Analog gilt auch $(-1) \cdot a = -a$.
				\item $(-1) \cdot (-1) = -(-1) = 1$
				\item 
				\begin{align*}
					(-a) \cdot (-b) 
					&= a \cdot (-1) \cdot (-b)\\
					&= a \cdot (-(-b))\\
					&= a \cdot b
				\end{align*}
			\end{enumerate}
		\end{proof}
		\ex{6.8} 
		\begin{enumerate}
			\item Jeder Ringhomomorphismus bildet Einheiten auf Einheiten ab.
			\item Jeder Ringhomomorphismus $K \to R$ von einem Körper $K$ in einen Ring $R$ mit $0_R \neq 1_R$ ist injektiv.
		\end{enumerate}
		\begin{proof}
			Sei $\phi : R \to S$ ein Ringhomomorphismus.
			\begin{enumerate}
				\item Sei $x$ eine Einheit. Dann gilt
				\begin{align*}
					1_S 
					&= \phi(1_R)\\
					&= \phi(x \cdot x^{-1})\\
					&= \phi(x) \cdot \phi(x^{-1})
				\end{align*}
				und analog $1_S  = \phi(x^{-1}) \cdot \phi(x)$. Somit ist $\phi(x)$ ebenfalls eine Einheit, mit $\phi(x)^{-1} = \phi(x^{-1})$.
				\item Sei nun $K$ ein Körper und $\phi : K \to R$ ein Ringhomomorphismus. Dann ist jedes $x \in K \setminus \lr{0_K}$ invertierbar, also gemäß der letzten Teilaufgabe ebenso jedes $\phi(x) \in \im~ \phi \setminus \lr{0_R}$. 
				\newpar
				Angenommen, $\phi$ wäre nicht injektiv. Dann gäbe es $x,y \in K$ mit $x - y \neq 0_K$, aber $\phi(x) - \phi(y) = 0_R$. Da $x - y \neq 0_K$, ist $(x-y)$ invertierbar. Es folgt:
				\begin{align*}
					\phi(x) - \phi(y) &= 0_R\\
					\implies \phi(x-y) &= 0_R\\
					\implies \phi(x-y)\phi(x-y)^{-1} &= 0_R\\
					\implies 1_R &= 0_R
				\end{align*}
				Ist nun $1_R \neq 0_R$ gegeben, ist dies ein Widerspruch, also $\phi$ injektiv. 
			\end{enumerate}
		\end{proof}
		\newpar
		\ex{7.9} $\bZ[X]$ ist kein Hauptidealring.
		\begin{proof}
			Wäre $\bZ[X]$ ein Hauptidealring, so könnte $\angles{2,X}$ von einem einzigen Element $a$ erzeugt werden. Dieses müsste ein gemeinsamer Teiler von $2$ und $X$ sein, also $a = \pm 1$. Es müsste also $\angles{\pm 1} = \angles{2,X}$. Allerdings gilt $1 \notin \angles{2,X}$.
		\end{proof}
		\newpar
		\ex{7.10} Im Quotientenring eines faktoriellen Rings $R$ nach einem Hauptideal zu einem Element $a \neq 0$ ist das Bild von $b \in \bR$ genau dann kürzbar, wenn $a$ und $b$ teilerfremd sind.
		\begin{proof}
			\begin{enumerate}
				\item Seien $a,b \in R$ gegeben, sodass $\ol b \in R / \angles{a}$ kürzbar ist. Wären $a$ und $b$ nicht teilerfremd, gäbe es ein irreduzibles $p \in R$ mit $p \mid a$ und $p \mid b$. Es gibt also ein $d \in R$, sodass $pd = a$. Es gilt außerdem $d \notin \angles{a}$, sonst wäre $p$ eine Einheit, also nicht irreduzibel. Nun ist gleichzeitig $b$ ein Vielfaches von $p$, also $bd$ ein Vielfaches von $pd = a$, also
				\begin{align*}
					\ol{b \cdot d} = \ol 0 = \ol {b \cdot 0}
				\end{align*}
				Mit $d \neq 0$. Also wäre $\ol b$ nicht kürzbar. Dies ist ein Widerspruch, also müssen $a$ und $b$ doch teilerfremd sein.
				\item 
			\end{enumerate}
		\end{proof}
		\newpar
		\ex{8.15} Ist $R$ ein faktorieller Ring mit Bruchkörper $K$ und sind $P,Q \in K[X]$ normierte Polynome mit $PQ \in R[X]$, so folgt bereits $P, Q \in R[X]$
		\newpar
		\ex{8.16} Sei $p$ prim und $r \geq 1$. So erfüllt das $p^r$-te Kreisteilungspolynom die Identität
		\begin{align*}
			\Phi_{p^r}(X) = \Phi_p\lr(X^{p^{r-1}})
		\end{align*}
	\clearpage
	\section*{\centering - Körper -}
		\ex{6.7} Man konstruiere einen Körper mit $169 = 13^2$ Elementen.
		\begin{proof}
			Gemäß der folgenden Aufgabe ist $\bF_{13}/\angles{X^2 - 7}$ ein solcher Körper.
		\end{proof}
		\newpar
		\ex{6.7.B} Sei $K$ ein endlicher Körper mit $q$ Elementen und $P$ vom Grad $n$ über $K$ irreduzibel. Dann ist $K/\angles{P}$ ein Körper mit $q^n$ Elementen.
		\newpar
		\ex{7.11} Sei $k$ ein Körper.
		\begin{enumerate}
			\item Alle Polynome vom Grad $1$ sind irreduzibel in $k[X]$.
			\item Ist $P \in k[X]$ irreduzibel und vom Grad $\deg P > 1$, so hat $P$ keine Nullstelle in $k$.
			\item Ist $P \in k[X] \setminus k$ vom Grad $\deg P \leq 3$ und hat $P$ keine Nullstelle in $k$, so ist $P$ irreduzibel in $k[X]$.
			\item Ist $k$ algebraisch abgeschlossen, so sind die irreduziblen Polynome in $k[X]$ genau die Polynome vom Grad $1$.
			\item Man gebe ein Polynom positiven Grades in $\bR[X]$ an, das keine Nullstelle hat, aber dennoch nicht irreduzibel ist.
		\end{enumerate}
		\newpar
		\ex{8.13} Die Definition der Summation in Bruchkörpern über Integritätsringen ist wohldefiniert.
		\newpar
		\ex{9.20} Sei $K \subset L$ ein Körper der Charakteristik ungleich zwei. Sei $\alpha, \beta \in L^\times$ mit $\alpha^2, \beta^2 \in K$: Dann gilt $K(\alpha) = K(\beta)$ genau dann, wenn $\alpha/\beta \in K$.
		\newpar
		\ex{10.21.b} Man bringe das Inverse von $7 + \sqrt[3]{2}$ in diese Standardform.
		\newpar
		\ex{10.23} Seien $R \supset K$ ein kommutativer Ring mit einem Teilring, der sogar ein Körper ist. Sei $R$ endlichdimensional als $K$-Vektorraum und ein Integritätsring. Dann ist $R$ auch selbst bereits ein Körper.
		\newpar
		\ex{11.25} Man zeige, dass es im Polynomring über einem endlichen Körper irreduzible Polynome von jedem positiven Grad gibt.
		\newpar
		\ex{11.26} Man gebe die Verknüpfungstafeln des Körpers mit vier Elementen an.
		\newpar
		\ex{12.1} Gegeben eine endliche Körpererweiterung $K \subset L$ ist jedes Polynom aus dem Polynomring $L[X]$ Teiler eines Polynoms aus $K[X]$
		\newpar
		\ex{12.2} Sei $p$ eine Primzahl und seien $\zeta, \xi \in \bC$ zwei primitive $p$-te Einheitswurzeln.
		\begin{enumerate}
			\item Es gilt $\bQ(\zeta) = \bQ(\xi)$.
			\item Es gibt genau einen Körperhomomorphismus $\bQ(\zeta) \to \bQ(\zeta)$ mit $\zeta \mapsto \xi$.
		\end{enumerate}
		\newpar
		\ex{12.3.a} Seien $M/L$ und $L/K$ algebraische Körpererweiterungen.
		\begin{enumerate}
			\item Ist $M/K$ normal, so ist auch $M/L$ normal.
			\item Sind $L_1 L_2 \subset M$ normale Körpererweiterungen von $K$, so ist $L_1 \cap L_2$ normal über $K$.
		\end{enumerate}
		\ex{12.3.b} Man finde Körper $M \supset L \supset K$, sodass $M/L$ und $L/K$ normal sind, aber $M/K$ nicht.
		\newpar
		\ex{12.4} Es sei $K$ ein Körper, $P \in K[X]$ ein Polynom vom Grad $n$ und $L/K$ der Zerfallungskörper von $P$. Man zeige:
		\begin{enumerate}
			\item $[L:K] \leq n!$
			\item $[L:K] \mid n!$
		\end{enumerate}
		\newpar
		\ex{13.2} Gegeben ein Körper $K$ und $p,q \in K$ ist das Polynom $P = X^3 + pX + q$ genau dann nicht separabel, wenn $27q^2 + 4p^3 = 0$.
		\begin{proof}
			Per Definition ist $P$ nicht separabel, wenn es in seinem Zerfallungskörper keine mehrfachen Nullstellen hat, also wenn $P$ und $P'$ über dem Zerfallungskörper keine gemeinsamen Nullstellen haben.
			\begin{enumerate}
				\item 
			Angenommen, es gäbe eine gemeinsame Nullstelle $\alpha$. Dann hätten wir:
			\begin{align*}
				P(\alpha) = \alpha^3 + p\alpha + q = 0
			\end{align*}
			und:
			\begin{align*}
				P'(\alpha) = 3\alpha^2 + p = 0
			\end{align*}
			\begin{enumerate}
				\item Angenommen, $\char~ K = 3$. Dann hätten wir
				\begin{align*}
					3\alpha^2 + p = 0\\ 
					\implies p = 0\\
					\implies 27q^2 + 4p^3 = 4p^3 = 0
				\end{align*}
				\item Angenommen, $\char K \neq 3$. Dann gilt:
				\begin{align*}
					 P'(\alpha) = 3\alpha^2 + p = 0 \implies \alpha^2 = -\frac{p}{3},
				\end{align*}
				also folgt:
				\begin{align*}
					0 &= P(\alpha)\\ 
					&= \alpha^3 + p\alpha + q\\
					&= \alpha \cdot \lr(-\frac{p}{3}) + p\alpha + q\\
					&= \frac{2p}{3}\alpha + q
				\end{align*}
				\begin{enumerate}
					\item In $\char~ K = 2$ erhalten wir $q = 0$, es folgt
					\begin{align*}
						27q^2 + 4p^3 = 27q^2 = 0
					\end{align*}
					\item
					In $\char K \notin \lr{2,3}$ stellen wir die Gleichung um zu:
					\begin{align*}
						&\frac{2p}{3}\alpha + q = 0\\
						\Leftrightarrow  &\frac{2p}{3}\alpha = -q\\
						\Leftrightarrow  &\alpha = -\frac{3q}{2p}\\
					\end{align*}
					Setzen wir dies in $P$ ein, erhalten wir:
					\begin{align*}
						\lr(-\frac{3q}{2p})^3 + \lr(-\frac{3q}{2p}) + q &= 0\\
					\Leftrightarrow -\frac{27q^3}{8p^3} -\frac{3q}{2p} + q &= 0\\
					\Leftrightarrow -27q^3 - 12qp^3 + 8qp^3 &= 0\\
					\Leftrightarrow 27q^3 + 4p^3 &= 0
					\end{align*}
				\end{enumerate}
			\end{enumerate} 
			\item Sei nun $27q^3 + 4p^3 = 0$. Dann folgt in Charakteristik $2$ $q = 0$ und in Charakteristik $3$ $p = 0$. Für $p = 0$ folgt $P = \alpha^3 + q = (X + \sqrt[3]{q})^3$, und für $q = 0$ folgt $P = X^3 + pX = X (X^2 + p)$, und in anderer Charakteristik kann wie in der vorherigen Rechnung $27q^3 + 4p^3 = 0$ auf $P\lr(-\frac{3q}{2p}) = 0$ zurückgeführt werden.
			\end{enumerate}
		\end{proof}
		\newpar
		\ex{13.3} Seien $M \supset L \supset K$ Körper. Ist $M/L$ separabel und $L/K$ separabel, so ist $M/K$ separabel.
		\newpar
		\ex{13.4} Seien $K$ ein Körper, $a \in K^\times$ und $n \geq 1$.
		\begin{enumerate}
			\item Im Zerfällungskörper des Polynoms $X^n - a$ zerfällt auch das Polynom $X^n - 1$ stets in Linearfaktoren.
			\item Im Zerfallungskörper des Polynoms $X^n-1$ zerfällt $X^n-a$ nicht notwendig in Linearfaktoren.
		\end{enumerate}
		\begin{proof}
		\begin{enumerate}
			\item Seien $\alpha_i$ die Nullstellen von $X^n - a$. Dann sind $\alpha_i/\alpha_1$ Nullstellen von $X^n - 1$.
			\item Über $\bQ$ zerfällt $X^2 - 1$ in Linearfaktoren, aber nicht $X^2 - 2$.
		\end{enumerate}
		\end{proof}
	\clearpage
	\section*{\centering - Zahlen -}
		\ex{0.1} Man führe an einem Beispiel den euklidischen Algorithmus aus und finde eine Darstellung des ggT als Linearkombination der beiden Ausgangszahlen.
		\newpar
		\ex{2.2} Man gebe alle Zahlen an, die bei Division durch $8$ Rest $4$ lassen, bei Division durch $13$ Rest $2$, und bei Division durch $11$ Rest $9$.
		\newpar
		\ex{3.1} \tbf{(Verallgemeinerte Fermat'sche Kongruenz)} Gegeben paarweise verschiedene Primzahlen $p_1, \hdots, p_r$ und eine Zahl $e$, sodass für alle $i \in \lr{1, \hdots r}$  $e \equiv 1 \mod (p_i-1)$ gilt, gilt für alle $a \in \bZ$ die Kongruenz
		\begin{align*}
			a^e \equiv a \mod (p_1 \cdot \hdots \cdot p_r)
		\end{align*}
		\newpar
		\ex{3.4} Gibt es ein Vielfaches von $23$, dessen letzte Ziffern $45$ sind? Gibt es eine Potenz von $23$, deren letzte Ziffern $45$ sind?
		\newpar
		\ex{6.5} Man finde das multiplikative Inverse der Nebenklasse von $21$ in $\bF_{31}$.
		\newpar
		\ex{6.6} Sei $n \in \bN$ mit $n \equiv 7 \mod 8$. Dann ist $n$ keine Summe von drei Quadraten.
		\newpar
		\ex{7.12} Man schreibe $130$ als Produkt von Gaußprimzahlen.
		\newpar
		\ex{9.19} Was ist die Summe der dritten Potenzen der vier Komplexen Nullstellen $\lambda_1, \lambda_2, \lambda_3, \lambda_4$ des Polynoms $X^4 + 3X^3 - 5X^2 + X + 13$?
		\newpar
		\ex{10.21} Alle Elemente von $\bQ(\sqrt[3]{2})$ lassen sich eindeutig in der Form
		\begin{align*}
			a + b \sqrt[3]{2} + c \sqrt[3]{2}^2
		\end{align*}
		mit $a,b,c \in \bQ$ schreiben.
		\newpar
		\ex{10.24.a} Ist $\sqrt 2 + \sqrt 3$ algebraisch über $\bQ$?
		\newpar
		\ex{10.24.b} Was ist das Minimalpolynom von $\sqrt 2 + \sqrt 3$ über $\bQ$?
		\newpar
		\ex{10.24.c} Liegt $\sqrt{2}$ in $\bQ(\sqrt{2} + \sqrt{3})$?
		\newpar
		\ex{11.27} Man zeige mithilfe der Axiomatischen Definition der natürlichen Zahlen, dass es für $a,b \in \bN$ mit $b \neq 0$ eindeutig bestimmte $c,d \in \bN$ gibt, sodass $d < b$ und $a = bc + d$.
		\newpar
		\ex{11.28} Man gebe einen Körperisomorphismus $\bF_7(\sqrt 3) \iso \bF_7(\sqrt 5)$ an.
	\clearpage
	\section*{\centering - Sonstiges zu Polynomen -}
		\ex{8.14} Gegeben ein Polynomen
		\begin{align*}
			a_n T^n + \hdots + a_1T + a_0
		\end{align*}
		mit ganzzahligen Koeffizienten und eine rationale Wurzel $p/q$ mit $p,q$ teilerfremden ganzen Zahlen ist $p$ ein Teiler von $a_0$ und $q$ ein Teiler von $a_n$
		\begin{proof}
			\begin{align*}
				\sum_{k = 0}^n a_k \lr(\frac{p}{q})^k &= 0\\
				\implies a_0 + \lr(\sum_{k = 1}^{n-1} a_k \lr(\frac{p}{q})^k) + a_n \lr(\frac{p}{q})^n &= 0\\
				\implies p \mid a_0 &= - \lr(\sum_{k = 1}^{n-1} a_k \lr(\frac{p}{q})^k) - a_n \lr(\frac{p}{q})^n\\
				q \mid a_n &= -\lr(\frac{q}{p})^n\lr(a_0 + \lr(\sum_{k = 1}^{n-1} a_k \lr(\frac{p}{q})^k))
			\end{align*}
		\end{proof}
		\newpar
		\ex{9.17} $P = X^7 - 9$ ist irreduzibel in $\bZ[X]$.
		\begin{proof}
			Sei $I : \bZ[X] \to \bZ[Y]$ die Injektion $X \to Y^2$. Dann gilt $I(P) = Y^{14} - 9 = (Y^7 - 3)(Y^7 + 3) := P_1P_2$. Nach Eisenstein sind $P_1$ und $P_2$ über $\bQ$ irreduzibel.
			\newpar
			Sei nun $P = QR$, sodass $R$ keine Einheit ist. Dann gilt $I(P) = I(QR) = I(Q) I(R)$, also $I(R) \mid I(P)$, also $I(R) \in \lr{P_1, P_2}$. Allerdings besteht das Bild von $I$ nur aus Polynomen mit geradem Grad, wir erhalten einen Widerspruch.
		\end{proof}
		\newpar
		\ex{9.18} Man faktorisiere $(X^n + Y^n)$ über $\bC[X,Y]$.
		\begin{proof}
			\begin{align*}
				X^n + Y^n &= \frac{X^{2n} + Y^{2n}}{X^n - Y^n}\\
				&= \frac{Y^{2n} \lr((\frac{X}{Y})^{2n} - 1)}{Y^n \lr((\frac{X}{Y})^{n} - 1)}\\
				&= \frac{Y^{2n} \prod_{\zeta^{2n} = 1}\lr(\frac{X}{Y} - \zeta)}{Y^n \prod_{\xi^n = 1}\lr(\frac{X}{Y} - \xi)}\\
				&= \frac{\prod_{\zeta^{2n} = 1}\lr(X - \zeta Y)}{\prod_{\xi^n = 1}\lr(X - \xi Y)}\\
				&= \prod_{\substack{\zeta^{2n} = 1, \\\zeta^n \neq 1}}\lr(X - \zeta Y)
			\end{align*}
		\end{proof}´
		\newpar
		\ex{10.22.B} Das Minimalpolynom von $z = a + bi \in \bC \setminus \bR$ über $\bR$ ist
		\begin{align*}
			P = (X - z)(X - \ol z)
		\end{align*}
		\begin{proof}
			Da $z \notin \bR$ kann es kein Minimalpolynom vom Grad $1$ geben. Gleichzeitig gilt:
			\begin{align*}
				&(X - z)(X - \ol z) = (X - (a + bi))(X - (a - bi))\\
				=& X^2 - X(a+bi) - X(a-bi) + (a+bi)(a-bi)\\
				=& X^2 - 2aX + a^2 + b^2 \in \bR
			\end{align*} 
			also ist $P \in \bR[X]$ normiert vom Grad $2$ mit $z$ als Nullstelle, also das Minimalpolynom.
		\end{proof}
		\newpar
		\ex{13.1} Man finde alle komplexen mehrfachen Nullstellen des Polynoms
		\begin{align*}
			P = X^4 - 4X^3 + 5X^2 - 4X + 4
		\end{align*}
		\begin{proof}
			$\lambda$ ist genau dann eine mehrfache Nullstelle von $P$, wenn es eine Nullstelle der Ableitung $P'$ ist. Die Ableitung ist:
			\begin{align*}
				P' = 4X^3 - 12X^2 + 10X - 4
			\end{align*}
			Wir berechnen den größten gemeinsamen Teiler der beiden Polynome, und den dazugehörigen Rest, durch Polynomdivision:
			\begin{align*}
				&X^4 - 4X^3 + 5X^2 - 4X + 4\\
				= &\lr(\frac{1}{4} X) \cdot \lr(4X^3 - 12X^2 + 10X - 4) - X^3 + 2.5X^2 - 3X + 4\\
				\\
				&-X^3 + 2.5X^2 - 3X + 4\\
				= &\lr(- \frac{1}{4}) \cdot \lr(4X^3 - 12X^2 + 10X - 4) - 0.5X^2 - 0.5X + 3
			\end{align*}
			\begin{align*}
				\implies P &= \lr(\frac{1}{4}X - \frac{1}{4}) P' + (- 0.5X^2 - 0.5X + 3)\\
				&:= t_1 \cdot P' + r_1
			\end{align*}
			Wir teilen nun $P'$ durch $r_1$ mit Rest:
			\begin{align*}
				&4X^3 - 12X^2 + 10X - 4\\ 
				= &\lr(-8X) \cdot \lr(- 0.5X^2 - 0.5X + 3) - 16X^2 + 34X - 4\\
				\\
				&-16X^2 + 34X - 4\\
				=&\lr(32) \cdot \lr(- 0.5X^2 - 0.5X + 3) + 50X - 100\\
			\end{align*}
			\begin{align*}
				P' 
				&= (-8X + 32) \cdot r_1 + 50X - 100\\
				&:= t_2 \cdot r_1 + r_2
			\end{align*}
			Zuletzt teilen wir $r_1$ durch $r_2$ ohne Rest:
			\begin{align*}
				(- 0.5X^2 - 0.5X + 3) : (50X - 100) &= -\frac{1}{100}X\\
				(-0.5X^2 + X)\\
				(-1.5X + 3): (50X - 100) &= -\frac{3}{100}\\
			\end{align*}
			Also ist $r_2 = 50X - 100 = 50 (X - 2)$ der größte gemeinsame Teiler von $P$ und $P'$. Insbesondere ist $2$ eine Nullstelle von $P'$, also eine mehrfache Nullstelle von $P$, also ist $(X-2)^2 = X^2 - 4X + 4$ ein Faktor von $P$. Wir erhalten durch erneute Polynomdivision:
			\begin{align*}
				&X^4 - 4X^3 + 5X^2 - 4X + 4 : X^2 - 4X + 4 = X^2 + 1\\
				&\underline{X^4 - 4X^3 + 4X^2}\\
				&\phantom{X^4 - X^3 + 4} X^2 - 4X + 4\\
				&\phantom{X^4 - X^3 + 4} \underline{X^2 - 4X + 4}\\
				&\phantom{X^4 - X^3 + 4 X^2 - 4X +}~0
			\end{align*}
			Also gilt $P = (X-2)^2(X-i)(X+i)$.
		\end{proof}
	\clearpage
	\section*{\centering - Sonstiges -}
		\ex{0.3} Seien $(M, \top)$ ein Monoid und $e$ sein neutrales Element. Man zeige: $M$ ist genau dann eine Gruppe, wenn es für jedes $a \in M$ ein $\ol a \in M$ gibt, sodass $\ol a \top a = e$, und dieses Element $\ol a$ ist dann notwendig bereits das Inverse von $a$ in $M$.
		\begin{proof}
			\begin{align*}
				  a \top \ol a 
				  &= e \top (a \top \ol a)\\
				  &= (\ol{\ol a} \top \ol a) \top a \top \ol a\\
				  &= \ol{\ol a} \top (\ol a \top a) \top \ol a\\
				  &= \ol{\ol a} \top e \top \ol a\\
				  &= \ol{\ol a} \top \ol a\\
				  &= e
			\end{align*}
			Also ist das Linksinverse auch das Rechtsinverse, also bildet $M$ eine Gruppe.
		\end{proof}
		\newpar
		\ex{0.4} Ein endliches Monoid $(M, \top)$, in dem für jedes Element $a \in M$ die Multiplikationsabbildung injektiv ist, muß bereits eine Gruppe sein.
		\begin{proof}
			Die Multiplikationsabbildung ist eine Abbildung $M \to M$, also folgt gemäß Schubfachprinzip aus Injektivität bereits Surjektivität. Somit existiert für jedes $a$ ein $b$ mit $a \top b = e$, und gemäß Übung 0.3 ist dieses $b$ auch das Linksinverse und $M$ ist eine Gruppe.
		\end{proof}
		\newpar
		\ex{1.1} Ein Element $a$ eines Monoids $M$ ist genau dann invertierbar, wenn es $b, c \in M$ gibt, sodass $b \top a = e = a \top c$.
		\begin{proof}
			\begin{align*}
				b 
				&= b \top e\\
				&= b \top (a \top c)\\
				&= (b \top a) \top c\\
				&= e \top c\\
				&= c
			\end{align*}
		\end{proof}
		\newpar
		\ex{4.2} Gegeben $p,q \in \bN$ mit $q \geq 1$ gibt es genau $\frac{(pq)!}{p!(q!)^p}$ Partitionen einer Menge mit $pq$ Elementen in $p$ Teilmengen mit jeweils $q$ Elementen.
\end{document}